% Ελληνική Περίληψη

\begin{abstract}\setlanguage{greek}
Ο έλεγχος νομικών συμβάσεων απαιτεί τον εντοπισμό συγκεκριμένων ρητρών μέσα σε εκτενή έγγραφα γραμμένα σε τυπική, τεχνική γλώσσα, γεγονός που καθιστά τη διαδικασία αργή και δαπανηρή για τις νομικές ομάδες. Τα μεγάλα γλωσσικά μοντέλα μπορούν να απαντούν σε ερωτήσεις πάνω σε κείμενο, όμως δεν μπορούν να απαντήσουν σε ερωτήσεις για μια συγκεκριμένη σύμβαση που δεν είδαν ποτέ κατά την εκπαίδευσή τους, και τείνουν να παράγουν απαντήσεις που διατυπώνονται με βεβαιότητα αλλά δεν τεκμηριώνονται από καμία πηγή. Η παρούσα διπλωματική εργασία υλοποιεί και αξιολογεί ένα συνομιλιακό σύστημα απάντησης ερωτήσεων για εμπορικές συμβάσεις, βασισμένο στην τεχνική της Ανάκτησης-Επαυξημένης Παραγωγής (Retrieval-Augmented Generation, RAG). Το σύστημα ανακτά τα σχετικά αποσπάσματα από μια σύμβαση και τα παρέχει ως πλαίσιο στο γλωσσικό μοντέλο, ώστε κάθε απάντηση να μπορεί να αναχθεί στο απόσπασμα του κειμένου από το οποίο προέκυψε. Παρέχεται ως διαδραστική διαδικτυακή εφαρμογή και εκτελείται εξ ολοκλήρου σε μία μόνο κάρτα γραφικών καταναλωτικού επιπέδου (GPU), χρησιμοποιώντας δύο τοπικά μοντέλα, το Llama 3.1 8B και το Mistral 7B, με κβαντισμό 4 bit. Η ανάκτηση συνδυάζει πυκνή σημασιολογική αναζήτηση με αναζήτηση λέξεων-κλειδιών BM25 και περιορίζει κάθε ερώτημα στη σύμβαση που κατονομάζεται, ενώ το σύστημα υποστηρίζει τρεις στρατηγικές τεμαχισμού και δύο στρατηγικές συνομιλιακής μνήμης. Με τη χρήση του συνόλου δεδομένων CUAD, το οποίο περιλαμβάνει 510 επισημειωμένες συμβάσεις, αξιολογήθηκαν έξι διαμορφώσεις με τέσσερις μετρικές RAGAS, δύο μετρικές τύπου LLM-as-a-judge και ένα σύστημα αναφοράς χωρίς ανάκτηση (baseline). Η ανάκτηση αύξησε την ορθότητα των απαντήσεων από 0,36 σε 0,77 για το ισχυρότερο μοντέλο, ενώ ο αναδρομικός ή ο σημασιολογικός τεμαχισμός με το Llama 3.1 8B έδωσε τις πιο ακριβείς απαντήσεις.

\vspace*{\fill}
\noindent{\large\bf{Λέξεις-κλειδιά:}}\\ 
Retrieval-Augmented Generation, Μεγάλα Γλωσσικά Μοντέλα, Έλεγχος Νομικών Συμβάσεων, Απάντηση Ερωτήσεων, CUAD, Υβριδική Ανάκτηση, Αξιολόγηση RAGAS
\end{abstract}







